\documentclass[11pt,a4paper]{article}

\usepackage[utf8]{inputenc}
\usepackage[T1]{fontenc}
\usepackage[spanish]{babel}
\usepackage{amsmath,amssymb,mathtools}
\usepackage{geometry}
\usepackage{lmodern}
\usepackage{microtype}

\geometry{margin=2.4cm}
\setlength{\parindent}{0pt}
\setlength{\parskip}{0.45em}

\newcommand{\R}{\mathbb{R}}

\begin{document}

\begin{center}
{\LARGE\bfseries Geometría Lineal}\\[0.4em]
{\Large Tema 6: Geometría proyectiva}\\[0.4em]
{\large Soluciones}
\end{center}

\section{Ejercicio 1}
\subsection{Solución}

Sean:

\[
P=(1:2:8),
\qquad
Q=(-7:0:3).
\]

La recta que pasa por ambos puntos viene dada por:

\[
\begin{vmatrix}
x_1&x_2&x_3\\
1&2&8\\
-7&0&3
\end{vmatrix}
=0.
\]

Desarrollando:

\[
\boxed{
6x_1-59x_2+14x_3=0.
}
\]

\section{Ejercicio 2}
\subsection{Solución}

Sean:

\[
P=(1:-1:2),
\qquad
Q=(2:1:1).
\]

La recta proyectiva está formada por los puntos:

\[
\lambda(1,-1,2)+\mu(2,1,1),
\qquad
(\lambda,\mu)\neq(0,0).
\]

Por tanto, sus ecuaciones paramétricas son:

\[
\boxed{
\begin{cases}
\alpha x_1=\lambda+2\mu,\\
\alpha x_2=-\lambda+\mu,\\
\alpha x_3=2\lambda+\mu,
\end{cases}
\qquad
(\alpha,\lambda,\mu)\neq(0,0,0).
}
\]

La ecuación implícita se obtiene mediante:

\[
\begin{vmatrix}
x_1&x_2&x_3\\
1&-1&2\\
2&1&1
\end{vmatrix}
=0.
\]

Así:

\[
\boxed{
x_1-x_2-x_3=0.
}
\]

\section{Ejercicio 3}
\subsection{Solución}

La recta \(r\) del ejercicio anterior tiene parametrización:

\[
(x_1:x_2:x_3)
=
(\lambda+2\mu:-\lambda+\mu:2\lambda+\mu).
\]

La recta:

\[
s:
x_1+x_2+x_3=0
\]

impone:

\[
2\lambda+4\mu=0.
\]

Por tanto:

\[
\lambda=-2\mu.
\]

Tomando:

\[
\mu=1,
\qquad
\lambda=-2,
\]

se obtiene:

\[
(0:3:-3).
\]

Por tanto:

\[
\boxed{
r\cap s=(0:1:-1).
}
\]

\section{Ejercicio 4}
\subsection{Solución}

\textbf{a)}

Las matrices de coeficientes de las tres rectas forman:

\[
M_a=
\begin{pmatrix}
1&-1&1\\
3&0&1\\
1&-1&-1
\end{pmatrix}.
\]

Se tiene:

\[
\det M_a=-6\neq0.
\]

Por tanto, las tres rectas no pasan por un mismo punto:

\[
\boxed{
\text{Las tres rectas forman un triángulo proyectivo.}
}
\]

\textbf{b)}

\[
M_b=
\begin{pmatrix}
1&1&1\\
1&3&-3\\
1&0&3
\end{pmatrix}.
\]

Se tiene:

\[
\det M_b=0.
\]

Resolviendo:

\[
M_b
\begin{pmatrix}
x_1\\x_2\\x_3
\end{pmatrix}
=0,
\]

se obtiene:

\[
(x_1:x_2:x_3)=(-3:2:1).
\]

Por tanto:

\[
\boxed{
\text{Las tres rectas son concurrentes en }(-3:2:1).
}
\]

\textbf{c)}

\[
M_c=
\begin{pmatrix}
1&1&0\\
1&0&1\\
0&1&1
\end{pmatrix}.
\]

Como:

\[
\det M_c=-2\neq0,
\]

se concluye:

\[
\boxed{
\text{Las tres rectas forman un triángulo proyectivo.}
}
\]

\section{Ejercicio 5}
\subsection{Solución}

Sean:

\[
P_1=(-1:0:0:1),
\]

\[
P_2=(0:1:1:1),
\]

\[
P_3=(1:2:a:b).
\]

Para que estén sobre una misma recta proyectiva debe cumplirse:

\[
P_3\in\mathcal L\{P_1,P_2\}.
\]

Por tanto:

\[
(1,2,a,b)
=
\lambda(-1,0,0,1)
+
\mu(0,1,1,1).
\]

Comparando coordenadas:

\[
-\lambda=1,
\]

\[
\mu=2,
\]

\[
a=\mu,
\]

\[
b=\lambda+\mu.
\]

Así:

\[
\boxed{
a=2,
\qquad
b=1.
}
\]

\section{Ejercicio 6}
\subsection{Solución}

La referencia es:

\[
R=\{P_1,P_2;P_3\},
\]

con:

\[
P_1=(1:1:0),
\qquad
P_2=(1:2:0),
\qquad
P_3=(0:1:0).
\]

Elegimos representantes normalizados:

\[
v_1=(-1,-1,0),
\]

\[
v_2=(1,2,0),
\]

de forma que:

\[
v_1+v_2=(0,1,0),
\]

que representa a \(P_3\).

Para:

\[
Q=(1:0:0),
\]

buscamos:

\[
(1,0,0)
=
\omega_1v_1+\omega_2v_2.
\]

Se obtiene:

\[
\omega_1=-2,
\qquad
\omega_2=-1.
\]

Por tanto:

\[
\boxed{
Q=(2:1)_R.
}
\]

\section{Ejercicio 7}
\subsection{Solución}

El sistema:

\[
R=\{P_1,P_2,P_3;P_4\}
\]

es la referencia canónica.

Para:

\[
R'=\{Q_1,Q_2,Q_3;Q_4\},
\]

se tiene:

\[
Q_1=(3:1:-3),
\]

\[
Q_2=(-1:0:5),
\]

\[
Q_3=(1:8:-1),
\]

\[
Q_4=(3:9:1).
\]

Además:

\[
Q_1+Q_2+Q_3=Q_4,
\]

por lo que la base ya está normalizada.

La matriz formada por los representantes de \(Q_1,Q_2,Q_3\) es:

\[
M=
\begin{pmatrix}
3&-1&1\\
1&0&8\\
-3&5&-1
\end{pmatrix}.
\]

Por tanto:

\[
X_R=M X_{R'}.
\]

Así:

\[
X_{R'}=M^{-1}X_R.
\]

Multiplicando por \(92\), que no altera las coordenadas homogéneas:

\[
\boxed{
\begin{pmatrix}
x_1'\\x_2'\\x_3'
\end{pmatrix}
\sim
\begin{pmatrix}
40&-4&8\\
23&0&23\\
-5&12&-1
\end{pmatrix}
\begin{pmatrix}
x_1\\x_2\\x_3
\end{pmatrix}.
}
\]

\section{Ejercicio 8}
\subsection{Solución}

La nueva referencia es:

\[
R'=\{P_4,P_2,P_3;P_1\}.
\]

Tomamos como base normalizada:

\[
v_1=(1,1,1),
\]

\[
v_2=(0,-1,0),
\]

\[
v_3=(0,0,-1),
\]

ya que:

\[
v_1+v_2+v_3=(1,0,0),
\]

que representa a \(P_1\).

La matriz de cambio es:

\[
M=
\begin{pmatrix}
1&0&0\\
1&-1&0\\
1&0&-1
\end{pmatrix}.
\]

Además:

\[
M^{-1}=M.
\]

Por tanto:

\[
\boxed{
\begin{pmatrix}
x_1'\\x_2'\\x_3'
\end{pmatrix}
=
\begin{pmatrix}
1&0&0\\
1&-1&0\\
1&0&-1
\end{pmatrix}
\begin{pmatrix}
x_1\\x_2\\x_3
\end{pmatrix}.
}
\]

Es decir:

\[
\boxed{
(x_1':x_2':x_3')
=
(x_1:x_1-x_2:x_1-x_3).
}
\]

\section{Ejercicio 9}
\subsection{Solución}

Para la referencia:

\[
R=
\{(1:1:0),(1:2:0),(0:1:1);(0:0:1)\},
\]

una base normalizada es:

\[
B_R=
\{(1,1,0),(-1,-2,0),(0,1,1)\}.
\]

Por tanto:

\[
M_R=
\begin{pmatrix}
1&-1&0\\
1&-2&1\\
0&0&1
\end{pmatrix}.
\]

Para:

\[
R'=
\{(1:0:0),(1:0:1),(0:1:1);(2:1:2)\},
\]

los tres primeros representantes ya están normalizados, ya que su suma es:

\[
(2,1,2).
\]

Así:

\[
M_{R'}=
\begin{pmatrix}
1&1&0\\
0&0&1\\
0&1&1
\end{pmatrix}.
\]

Para:

\[
P=(2:1:2),
\]

se obtiene:

\[
M_R^{-1}
\begin{pmatrix}
2\\1\\2
\end{pmatrix}
=
\begin{pmatrix}
5\\3\\2
\end{pmatrix}.
\]

Por tanto:

\[
\boxed{
P=(5:3:2)_R.
}
\]

Como \(P\) es el cuarto punto de \(R'\):

\[
\boxed{
P=(1:1:1)_{R'}.
}
\]

El cambio de \(R\) a \(R'\) es:

\[
M_{R'}^{-1}M_R
=
\begin{pmatrix}
2&-3&0\\
-1&2&0\\
1&-2&1
\end{pmatrix}.
\]

Así:

\[
\boxed{
X_{R'}
=
\begin{pmatrix}
2&-3&0\\
-1&2&0\\
1&-2&1
\end{pmatrix}
X_R.
}
\]

El cambio inverso es:

\[
\boxed{
X_R
=
\begin{pmatrix}
2&3&0\\
1&2&0\\
0&1&1
\end{pmatrix}
X_{R'}.
}
\]

\section{Ejercicio 10}
\subsection{Solución}

Sean:

\[
A=(1:-1),
\quad
B=(3:-1),
\quad
C=(1:0),
\quad
D=(2:-1).
\]

La razón doble es:

\[
[A,B,C,D]
=
\frac{
\begin{vmatrix}
1&-1\\
1&0
\end{vmatrix}
\begin{vmatrix}
3&-1\\
2&-1
\end{vmatrix}
}{
\begin{vmatrix}
1&-1\\
2&-1
\end{vmatrix}
\begin{vmatrix}
3&-1\\
1&0
\end{vmatrix}
}.
\]

Por tanto:

\[
\boxed{
[A,B,C,D]=-1.
}
\]

\section{Ejercicio 11}
\subsection{Solución}

Sean:

\[
A=(1:0:0),
\]

\[
B=(1:1:1),
\]

\[
C=(0:1:1),
\]

\[
D=(-2:1:1).
\]

\textbf{a)}

Tomando \(A\) y \(B\) como generadores de la recta:

\[
C=-A+B,
\]

\[
D=-3A+B.
\]

Por tanto:

\[
\boxed{
A,B,C,D\text{ están alineados}.
}
\]

\textbf{b)}

Respecto de la base \(\{A,B\}\):

\[
A=(1:0),
\qquad
B=(0:1),
\qquad
C=(-1:1),
\qquad
D=(-3:1).
\]

Así:

\[
\boxed{
[A,B,C,D]=3.
}
\]

\textbf{c)}

Sea:

\[
A'=(a:b).
\]

Imponemos:

\[
[A',B,C,D]=-1.
\]

Se obtiene:

\[
2a+3b=0.
\]

Podemos tomar:

\[
(a,b)=(-3,2).
\]

Por tanto:

\[
A'=-3A+2B.
\]

Así:

\[
\boxed{
A'=(-1:2:2).
}
\]

\section{Ejercicio 12}
\subsection{Solución}

Tomamos:

\[
A=(1:1:-1:1),
\]

\[
B=(0:1:1:-1).
\]

Los otros puntos se expresan como:

\[
C=2A-B,
\]

\[
D=4A-B.
\]

Por tanto, respecto de la referencia inducida en la recta:

\[
A=(1:0),
\]

\[
B=(0:1),
\]

\[
C=(2:-1),
\]

\[
D=(4:-1).
\]

Así:

\[
\boxed{
[A,B,C,D]=2.
}
\]

\section{Ejercicio 13}
\subsection{Solución}

Sea la matriz de la proyectividad:

\[
A=
\begin{pmatrix}
a&b\\
c&d
\end{pmatrix}.
\]

Como:

\[
f(1:0)=(2:-3),
\]

la primera columna es proporcional a:

\[
(2,-3).
\]

Como:

\[
f(0:1)=(1:1),
\]

la segunda columna es proporcional a:

\[
(1,1).
\]

Por tanto:

\[
A=
\begin{pmatrix}
2\alpha&\beta\\
-3\alpha&\beta
\end{pmatrix}.
\]

La condición:

\[
f(1:1)=(0:5)
\]

impone:

\[
2\alpha+\beta=0.
\]

Tomando:

\[
\alpha=1,
\qquad
\beta=-2,
\]

obtenemos:

\[
A=
\begin{pmatrix}
2&-2\\
-3&-2
\end{pmatrix}.
\]

Entonces:

\[
A
\begin{pmatrix}
7\\21
\end{pmatrix}
=
\begin{pmatrix}
-28\\-63
\end{pmatrix}.
\]

Por tanto:

\[
\boxed{
f(7:21)=(4:9).
}
\]

\section{Ejercicio 14}
\subsection{Solución}

La proyectividad verifica:

\[
f(1:0:0)=(0:0:1),
\]

\[
f(0:1:0)=(0:1:1),
\]

\[
f(1:1:1)=(1:0:1),
\]

\[
f(2:0:1)=(-1:2:3).
\]

Tomando las columnas de la matriz asociada y ajustando los escalares mediante las dos últimas condiciones, se obtiene:

\[
A=
\begin{pmatrix}
0&0&-1\\
0&-2&2\\
2&-2&-1
\end{pmatrix}.
\]

Por tanto:

\[
\boxed{
\begin{pmatrix}
x_1'\\x_2'\\x_3'
\end{pmatrix}
=
\begin{pmatrix}
0&0&-1\\
0&-2&2\\
2&-2&-1
\end{pmatrix}
\begin{pmatrix}
x_1\\x_2\\x_3
\end{pmatrix}.
}
\]

\section{Ejercicio 15}
\subsection{Solución}

La matriz de la homografía es:

\[
A=
\begin{pmatrix}
2&1&-1\\
0&2&-1\\
-3&-2&3
\end{pmatrix}.
\]

Los puntos fijos corresponden a autovectores de \(A\).

El polinomio característico es:

\[
(\lambda-1)^2(\lambda-5).
\]

Para:

\[
\lambda=1,
\]

se obtiene:

\[
V_1=\mathcal L\{(0,1,1)\}.
\]

Para:

\[
\lambda=5,
\]

se obtiene:

\[
V_5=\mathcal L\{(-4,-3,9)\}.
\]

Por tanto:

\[
\boxed{
P_1=(0:1:1),
\qquad
P_2=(-4:-3:9).
}
\]

\section{Ejercicio 16}
\subsection{Solución}

Los puntos son:

\[
A=(0:0:1),
\quad
B=(0:1:0),
\quad
C=(1:0:0),
\quad
D=(1:1:1).
\]

Buscamos una homografía que verifique:

\[
A\mapsto B,
\quad
B\mapsto C,
\quad
C\mapsto D,
\quad
D\mapsto A.
\]

Una matriz asociada es:

\[
\boxed{
A_f=
\begin{pmatrix}
1&-1&0\\
1&0&-1\\
1&0&0
\end{pmatrix}.
}
\]

Su polinomio característico es:

\[
(\lambda-1)(\lambda^2+1).
\]

El único autovalor real es:

\[
\lambda=1,
\]

con autovector:

\[
(1,0,1).
\]

Por tanto:

\[
\boxed{
P=(1:0:1)
}
\]

es el único punto fijo real.

\section{Ejercicio 17}
\subsection{Solución}

La matriz de la homografía es:

\[
A=
\begin{pmatrix}
0&-1&1\\
\beta&0&-1\\
1&-7&0
\end{pmatrix}.
\]

\textbf{a)}

El punto:

\[
P=(1:0:1)
\]

debe ser fijo.

Se tiene:

\[
A
\begin{pmatrix}
1\\0\\1
\end{pmatrix}
=
\begin{pmatrix}
1\\
\beta-1\\
1
\end{pmatrix}.
\]

Por tanto:

\[
\boxed{
\beta=1.
}
\]

Para este valor, los autovalores son:

\[
-3,\qquad1,\qquad2.
\]

Los puntos fijos son:

\[
\boxed{
(-1:2:5),
\qquad
(1:0:1),
\qquad
(3:-1:5).
}
\]

\textbf{b)}

Las rectas invariantes corresponden a autovectores de \(A^t\).

Se obtiene:

\[
\boxed{
-x_1+2x_2+x_3=0,
}
\]

\[
\boxed{
-3x_1-4x_2+x_3=0,
}
\]

\[
\boxed{
-x_1-3x_2+x_3=0.
}
\]

\section{Ejercicio 18}
\subsection{Solución}

La matriz es:

\[
A=
\begin{pmatrix}
2&1&-1\\
0&2&-1\\
-3&-2&3
\end{pmatrix}.
\]

Los puntos fijos son los autovectores de \(A\):

\[
\boxed{
(0:1:1),
\qquad
(-4:-3:9).
}
\]

Las rectas invariantes se obtienen a partir de los autovectores de \(A^t\).

Para:

\[
\lambda=1,
\]

se obtiene:

\[
(3,-1,1).
\]

Para:

\[
\lambda=5,
\]

se obtiene:

\[
(-1,-1,1).
\]

Por tanto, las rectas invariantes son:

\[
\boxed{
3x_1-x_2+x_3=0
}
\]

y:

\[
\boxed{
-x_1-x_2+x_3=0.
}
\]

\section{Ejercicio 19}
\subsection{Solución}

La matriz es:

\[
A=
\begin{pmatrix}
2&0&0\\
1&2&-1\\
1&0&2
\end{pmatrix}.
\]

Su único autovalor es:

\[
\lambda=2.
\]

El subespacio propio es:

\[
V_2=\mathcal L\{(0,1,0)\}.
\]

Por tanto, el único punto fijo es:

\[
\boxed{
P=(0:1:0).
}
\]

Para \(A^t\), el subespacio propio es:

\[
\mathcal L\{(1,0,0)\}.
\]

Así, la única recta invariante es:

\[
\boxed{
x_1=0.
}
\]

\section{Ejercicio 20}
\subsection{Solución}

La matriz de la homografía es:

\[
A=
\begin{pmatrix}
-11&-8&-16\\
16&11&20\\
1&1&3
\end{pmatrix}.
\]

El polinomio característico es:

\[
(\lambda-1)^3.
\]

El subespacio propio de \(A\) es:

\[
V_1=
\mathcal L\{(0,-2,1)\}.
\]

Por tanto, el único punto fijo es:

\[
\boxed{
P=(0:-2:1).
}
\]

Para la matriz traspuesta:

\[
A^t,
\]

el subespacio propio es:

\[
\mathcal L\{(3,2,4)\}.
\]

Por tanto, la única recta invariante es:

\[
\boxed{
3x_1+2x_2+4x_3=0.
}
\]

\end{document}
